% !TEX program = lualatex
% !TEX encoding = UTF-8 Unicode
\documentclass[11pt]{beamer}
\usepackage{luatexja}
\usepackage[haranoaji,deluxe]{luatexja-preset}
\usepackage{amsmath}
\usepackage{amsthm}
\setbeamertemplate{proof begin}{\noindent\textbf{証明}\quad}
\setbeamertemplate{proof end}{\qed \par}
\setbeamertemplate{bibliography item}[text]

\theoremstyle{plain}
\newtheorem{thm}{定理}
\newtheorem*{thm*}{定理}
\newtheorem{cor}{系}
\newtheorem{lem}{補題}
\newtheorem{assumption}{仮定}

\theoremstyle{definition}
\newtheorem{dfn}{定義}
\def\qedsymbol{}
% ------------------------------
%  pLaTeX で日本語を扱う設定例
% ------------------------------
% hyperref で日本語 PDF しおり等が化けないように

% 数式や画像，色の設定
\usepackage{amsmath,amssymb}
\usepackage{graphicx}
\usepackage{xcolor}

% ------------------------------
% Beamer のテーマや色の設定
% ------------------------------
\usetheme{Madrid}       % お好みのテーマに変更可能
\usecolortheme{beaver}  % 赤系配色の例

% スライド下部にページ番号を表示したい場合
\setbeamertemplate{footline}[frame number]
% ナビゲーションアイコンを消す場合
\setbeamertemplate{navigation symbols}{}


% for hyperref
\hypersetup{
	colorlinks=true, % リンクに色をつけない設定
	bookmarks=true, % 以下ブックマークに関する設定
	bookmarksnumbered=true,
	pdfborder={0 0 0},
	bookmarkstype=toc
}

% ------------------------------
% タイトル情報
% ------------------------------
%1
\title{データの内在次元を考慮したWasserstein距離での拡散モデルの収束解析}
%\subtitle{}
\author{枇榔 一真}
\institute{関根$\cdot$ 深澤 $\cdot$ 矢野 研究室}
\date{2026年7月9日}

% ------------------------------
\begin{document}

% ---- タイトルスライド ----
\begin{frame}
  \titlepage
\end{frame}
 %2
  \begin{frame}{目次}
    \begin{itemize}
      \item イントロダクション
      \begin{itemize}
        \item 拡散モデルの収束解析について
        \item 拡散モデルの収束解析について (KL ダイバージェンス)
        \item 拡散モデルの収束解析について (2-Wasserstein 距離)
        \item 2つの評価指標の比較
        \item データの内在次元について
        \item 内在次元を用いた拡散モデルの収束解析について
      \end{itemize}
      \item 設定
      \item 主定理の証明概要
    \end{itemize}
  \end{frame}
  \begin{frame}{拡散モデルの収束解析について}
    \begin{itemize}
      \item 拡散モデルは, 未知の確率測度 $p_{data}$ に従って独立に得られるデータ $x_1,\cdots, x_n$ を用いて $p_{data}$ に従うデータを新たに近似的に生み出す手法である.
      \item 近似的に生み出すデータが従う確率測度を $\hat{p}$ とした時に,何かしらの確率測度の評価指標 $\rho$ を用いて $\rho(p_{data},\hat{p})$ を評価することが収束解析の主な目的である.
      \item 先行研究では, 評価指標 $\rho$ として主にKLダイバージェンスと2-Wasserstein距離が用いられている.(KLダイバージェンスは距離の公理は満たさないがよく用いられている)
    \end{itemize}
  
  \end{frame}
  %-----------3-------------
  \begin{frame}{拡散モデルの収束解析について(KLダイバージェンス)}
    \begin{itemize}
      \item KLでの収束解析は \cite{benton_et_al_2024, chen_lee_lu_2023, chen_et_al_2023, koike_2026_follmer} など多くの先行研究がある. 
      \item データは $\mathbb{R}^d$ 上に分布するとしたとき, \cite{benton_et_al_2024} では $KL(p_{data},\hat{p})\leqq O(d)$ の評価が得られており, これは最適なオーダーであることが示されている.
      \item この結果は, $p_{data}$ に関する一般的な仮定のもとで得られている.
    \end{itemize}
  
  \end{frame}
  %------------4-----------
  \begin{frame}{拡散モデルの収束解析について(2-Wasserstein距離)}
    \begin{itemize}
        \item 2-Wasserstein距離での収束解析も近年増加している\cite{arsenyan_vardanyan_dalalyan_2025, beyler_bach_2025, bruno_sabanis_2025, gao_nguyen_zhu_2025} .
        \item しかし, これらの先行研究ではデータ分布 $p_{data}$ へのlog-concave性など, 一般的ではない仮定が置かれており, KLと比較して発展途上である.
    \end{itemize}
  \end{frame}

  %---------5-------------
  \begin{frame}{2つの評価指標の比較}
    \begin{itemize}
        \item KLダイバージェンス
        \begin{itemize}
          \item ギルザノフの定理を使うことで容易に評価が可能であり, 一般的な仮定のもとで最適と思われるオーダー評価ができている.
          \item KLダイバージェンスが近くても平均や共分散が近いとは限らないなどの問題点がある \cite{beyler_bach_2025}.
        \end{itemize}
        \item 2-Wasserstein距離
        \begin{itemize}
          \item 画像生成の評価指標として広く用いられているFréchet Inception Distance \cite{heusel_et_al_2017} が2-Wasserstein距離と関連があるため, 重要と言われている \cite{arsenyan_vardanyan_dalalyan_2025}.
        \end{itemize}
    \end{itemize}
  \end{frame}
  %---------6-------------
  \begin{frame}{データの内在次元について}
    \begin{itemize}
        \item 現実世界のデータ分布はごく少数の変数で記述されると考えられている. これは多様体仮説と呼ばれ, 機械学習の成功の基盤と考えられている \cite{pope_et_al_2021}.
        \item データ分布を記述する上で必要な変数の数を, データ分布の\textbf{内在次元 (intrinsic dimension)} と呼ぶ.
        \item \cite{pope_et_al_2021} は, MNIST、CIFAR-10、ImageNet などの代表的な画像データセットの内在次元推定を行った.
        \item 次元推定の結果, ImageNetでは各画像が $224\times 224 \times 3 = 15028$ の画素値を持つが, 内在次元は $26\sim 43$ 程度であるとされている.
    \end{itemize}
  \end{frame}
  %---------7-------------
    \begin{frame}{内在次元を用いた拡散モデルの収束解析について}
    \begin{itemize}
        \item データの内在次元を用いた拡散モデルの収束解析は \cite{huang_wei_chen_2026,azangulov_deligiannidis_rousseau_2024,li_yan_2024_low_dimensional} などによって行われており, 
        \cite{huang_wei_chen_2026} では $KL(p_{data},\hat{p})\leqq O(k)$ の評価が得られている.
        \item 上記の結果はKLダイバージェンスでの評価であり, 2-Wasserstein距離での収束解析は \cite{huang_wei_chen_2026} ではFuture Work とされている.
        \item 今回, \cite{huang_wei_chen_2026} の評価方法を参考に2-Wasserstein距離での評価を行った.
    \end{itemize}
  \end{frame}
  %---------8-------------
    \begin{frame}{拡散モデルの更新則}
      $X_0\sim p_{data}$, $W \sim \mathcal{N}(0,I_d)$ として, 確率過程 $\{X_t\}_{t\in[0,T]}$ を次を満たすものとして定める.
      \begin{equation}
         X_t\stackrel{\mathrm{d}}{=}\sqrt{1-\sigma^2_t}X_0+\sigma_tW \quad (\sigma_t=\sqrt{1-e^{-2t}}, \quad t\in[0,T])
      \end{equation}
      $X_t$ の密度関数を $p_t$ とし, $s_t(x):=\nabla\log{p_t(x)}$ の推定量を $\hat{s}_t(x)$ として, データを生成する過程 $\{Y_n\}_{n=0,1,\dots, N}$ を次のように定義する.
      \begin{align}
          &Y_0\sim \mathcal{N}(0,I_d), \notag\\
          &Y_{n+1}=\frac{1}{\sqrt{\alpha_n}}(Y_{n}+(1-\alpha_n)\hat{s}_{T-t_n}(Y_{n}))+\sqrt{\frac{(1-\alpha_n)(1-\bar{\alpha}_{n+1})}{1-\bar{\alpha_n}}}Z_n
      \end{align}
      ここで, $0=t_0<t_1<\dots < t_N<t_{N+1}=T$, $\alpha_n=e^{-2(t_{n+1}-t_n)}$ , $\bar{\alpha}_n=\prod\limits_{i=n}^N \alpha_i$ , $Z_n\overset{iid}{\sim}\mathcal{N}(0,I_d)$ である. 係数 $\alpha_n$ の値は \cite{ho_jain_abbeel_2020} で提案された.
    \end{frame}
    %---------9-------------
    \begin{frame}{SDEを通じた解釈}
      前頁のアルゴリズムの時間を連続化することでSDEを使った解釈が可能である \cite{song_et_al_2021} .
      \begin{equation}
          dX_t=-X_tdt+\sqrt{2}dW_t, \quad X_0\sim p_{data}, \quad t\in[0,T]
      \end{equation}
      の解 $\{X_t\}_{t\in[0,T]}$ に対して $X_t$ の密度関数を $p_t$ , $s_t(x):=\nabla\log{p_t(x)}$ として, 次のSDEを考える.
      \begin{equation}
          dY_t=(Y_t+2s_{T-t}(Y_t))dt+\sqrt{2}dB_t \quad t\in[0,T]
      \end{equation}
      $Y_0\stackrel{\mathrm{d}}{=} X_T$ としたとき, (4) の解を $\{Y_t\}_{t\in[0,T]}$ とすると, 任意の時刻 $t\in [0,T]$ において
      $Y_{T-t}\stackrel{\mathrm{d}}{=}X_t$ が成立する\cite{anderson82}.
    \end{frame}
    %---------10-------------
    \begin{frame}{Tweedieの公式}
      $\{X_t\}_{t\in[0,T]}$ を前頁のSDEの解として, 
      \begin{equation}
          \mu_t(x)=\mathbb{E}[X_0|X_t=x], \quad \Sigma_t(x)=Cov[X_0|X_t=x]
      \end{equation}
      とする. このとき, Tweedieの公式 \cite{efron_2011_tweedie} より, 
      \begin{equation}
          \mu_t(X_t)=\frac{1}{\sqrt{1-\sigma^2_t}}(X_t+\sigma^2_t\nabla \log{p_t(X_t)})=\frac{1}{\sqrt{1-\sigma^2_t}}(X_t+\sigma^2_ts_t(X_t))
      \end{equation}
      が得られる.
      これを用いると, 前頁の時間を反転したSDEは次のように書ける.
      \begin{align}
          dY_t&=(Y_t+\frac{2\sqrt{1-\sigma^2_{T-t}}}{\sigma^2_{T-t}}\mu_{T-t}(Y_t)-\frac{2}{\sigma^2_{T-t}}Y_t)dt+\sqrt{2}dB_t\\
          &=\{(1-\frac{2}{\sigma^2_{T-t}})Y_t+\frac{2\sqrt{1-\sigma^2_{T-t}}}{\sigma^2_{T-t}}\mu_{T-t}(Y_t)\}dt+\sqrt{2}dB_t
      \end{align}
    \end{frame}
    %---------11-------------
    \begin{frame}{実際の生成アルゴリズムとの対応}
      前頁の(8)式のSDE
      \begin{align}
          dY_t=\{(1-\frac{2}{\sigma^2_{T-t}})Y_t+\frac{2\sqrt{1-\sigma^2_{T-t}}}{\sigma^2_{T-t}}\mu_{T-t}(Y_t)\}dt+\sqrt{2}dB_t
      \end{align}
      に対して, スコア関数の推定量を$\hat{s}$ , $\hat{\mu}_t(x)=\frac{1}{\sqrt{1-\sigma^2_t}}(x+\sigma^2_t\hat{s}_t(x))$ として, $\{\hat{Y}_t\}_{t\in[0,T]}$ を, $\hat{Y}_0\sim \mathcal{N}(0,I_d)$ であるような次のSDEの解とする.
      \begin{align}
          d\hat{Y}_t=\{(1-\frac{2}{\sigma^2_{T-t}})\hat{Y}_t+\frac{2\sqrt{1-\sigma^2_{T-t}}}{\sigma^2_{T-t}}\hat{\mu}_{T-t_n}(\hat{Y}_{t_n})\}dt+\sqrt{2}dB_t \notag\\
          (t\in (t_n,t_{n+1}])
      \end{align}
    \end{frame}
    %---------12-------------
    \begin{frame}{実際の生成アルゴリズムとの対応}
      前頁の $\{\hat{Y}_t\}_{t\in[0,T]}$  に対して, 次の定理が成立する.
      \begin{theorem}
          任意の $n=0,\cdots, N-1$ に対して, ある独立同分布で標準正規分布に従う $\{Z_n\}_{n=0,\cdots, N-1}$ がが存在して, 次の関係が成り立つ.
          \begin{equation}
            \hat{Y}_{n+1}=\frac{1}{\sqrt{\alpha_n}}(\hat{Y}_{n}+(1-\alpha_n)\hat{s}_{T-t_n}(\hat{Y}_{n}))+\sqrt{\frac{(1-\alpha_n)(1-\bar{\alpha}_{n+1})}{1-\bar{\alpha_n}}}Z_n
          \end{equation}
      \end{theorem}
      この定理により, (10) のSDEは (2) のアルゴリズムを連続時間補完したものと解釈できることがわかる.
    \end{frame}
    %---------13-------------
    \begin{frame}{データの内在次元の仮定について}
      データの内在次元を考えるために, 集合の複雑さに関する概念を定義する.
      \begin{definition}
        $A\subset \mathbb{R}^d$ に対する, スケール $\epsilon_0$ における \textbf{covering number} $N^{cover}(A,\|\|_2,\epsilon_0)$ という値を, 次の関係を満たす $x_1,\cdots, x_s \in \mathbb{R}^d$ 
        が存在するような最小の自然数 $s$ として定義する.
        \begin{equation}
          A\subset \cup_{i=1}^sB(x_i,\epsilon_0)
        \end{equation}
        また, $\log{N^{cover}(A,\|\|_2,\epsilon_0)}$ を \textbf{metric entropy} という.
      \end{definition}
    \end{frame}

    \begin{frame}[allowframebreaks]{参考文献}
    \begin{thebibliography}{99}

      \bibitem{benton_et_al_2024}
      Benton, J., De Bortoli, V., Doucet, A. and Deligiannidis, G. (2024).
      Nearly $d$-linear convergence bounds for diffusion models via stochastic localization.
      In \textit{International Conference on Learning Representations}.
      \url{https://proceedings.iclr.cc/paper_files/paper/2024/hash/9edce07129d2ec4c511a8783eb3d81b5-Abstract-Conference.html}

      \bibitem{huang_wei_chen_2026}
      Huang, Z., Wei, Y. and Chen, Y. (2024; revised 2026).
      Denoising diffusion probabilistic models are optimally adaptive to unknown low dimensionality.
      \url{https://arxiv.org/abs/2410.18784}

      
      \bibitem{chen_lee_lu_2023}
      Chen, H., Lee, H. and Lu, J. (2023).
      Improved analysis of score-based generative modeling: user-friendly bounds under minimal smoothness assumptions.
      In \textit{Proceedings of the 40th International Conference on Machine Learning}.
      \url{https://proceedings.mlr.press/v202/chen23q.html}
    
    \bibitem{anderson82}
    Anderson, B. D. O. (1982).
    Reverse-time diffusion equation models.
    Stochastic Process. Appl. 12 313–326.
    \url{https://mathscinet.ams.org/mathscinet/article?mr=656280}
    
    \bibitem{haussmann_pardoux86}
    Haussmann, U. G. and Pardoux, E. (1986).
    Time reversal of diffusions.
    Ann. Probab. 14 1188–1205.
    \url{https://mathscinet.ams.org/mathscinet/article?mr=866342}
    
    \bibitem{tang_zhao_2024}
    Wenpin Tang and Hanyang Zhao (2024).
    Score-based Diffusion Models via Stochastic Differential Equations -- a Technical Tutorial.
    \url{https://arxiv.org/abs/2402.07487}

    \bibitem{efron_2011_tweedie}
    Efron, B. (2011).
    Tweedie's formula and selection bias.
    \textit{Journal of the American Statistical Association}
    106(496), 1602--1614.
    \url{https://doi.org/10.1198/jasa.2011.tm11181}
    
    \bibitem{chen_et_al_2023}
    Chen, S., Chewi, S., Li, J., Li, Y., Salim, A. and Zhang, A. R. (2023).
    Sampling is as easy as learning the score: theory for diffusion models with minimal data assumptions.
    \url{https://arxiv.org/abs/2209.11215}
    \bibitem{borgi_a_2019}
    Borji, A. (2019). Pros and cons of GAN evaluation measures. Comput.
    Vis. Image Underst. 179 41–65.
    \url{https://arxiv.org/abs/1802.03446}

    \bibitem{koike_2026_follmer}
    Koike, Y. (2026).
    A note on connections between the F\"ollmer process and the
    denoising diffusion probabilistic model.
    arXiv preprint arXiv:2605.18040.
    \url{https://arxiv.org/abs/2605.18040}

    \bibitem{arsenyan_vardanyan_dalalyan_2025}
    Arsenyan, V., Vardanyan, E. and Dalalyan, A. S. (2025).
    Assessing the quality of denoising diffusion models in Wasserstein distance:
    noisy score and optimal bounds.
    In \textit{Advances in Neural Information Processing Systems 38}.
    \url{https://proceedings.neurips.cc/paper_files/paper/2025/hash/1c26c389d60ec419fd24b5fee5b35796-Abstract-Conference.html}
    
    \bibitem{beyler_bach_2025}
    Beyler, E. and Bach, F. (2025).
    Convergence of deterministic and stochastic diffusion-model samplers:
    a simple analysis in Wasserstein distance.
    arXiv preprint arXiv:2508.03210.
    \url{https://arxiv.org/abs/2508.03210}
    
    \bibitem{bruno_sabanis_2025}
    Bruno, S. and Sabanis, S. (2025).
    Wasserstein convergence of score-based generative models under
    semiconvexity and discontinuous gradients.
    \textit{Transactions on Machine Learning Research}.
    \url{https://openreview.net/forum?id=vS9iVRB7XF}
    
    \bibitem{gao_nguyen_zhu_2025}
    Gao, X., Nguyen, H. M. and Zhu, L. (2025).
    Wasserstein convergence guarantees for a general class of
    score-based generative models.
    \textit{Journal of Machine Learning Research} 26(43), 1--54.
    \url{https://jmlr.org/papers/v26/24-0902.html}
    
    \bibitem{gentiloni_silveri_ocello_2025}
    Gentiloni Silveri, M. and Ocello, A. (2025).
    Beyond log-concavity and score regularity:
    improved convergence bounds for score-based generative models
    in $W_2$-distance.
    In \textit{Proceedings of the 42nd International Conference on Machine Learning},
    Proceedings of Machine Learning Research 267, 55631--55656.
    \url{https://proceedings.mlr.press/v267/silveri25a.html}
    
    \bibitem{pfarr_timofte_werner_2026}
    Pfarr, E., Timofte, R. and Werner, F. (2026).
    Analyzing the error of generative diffusion models:
    from Euler--Maruyama to higher-order schemes.
    arXiv preprint arXiv:2601.18425.
    \url{https://arxiv.org/abs/2601.18425}
    
    \bibitem{strasman_et_al_2025}
    Strasman, S., Ocello, A., Boyer, C., Le Corff, S. and Lemaire, V. (2025).
    An analysis of the noise schedule for score-based generative models.
    \textit{Transactions on Machine Learning Research}.
    \url{https://openreview.net/forum?id=BlYIPa0Fx1}
    
    \bibitem{yu_yu_2025}
    Yu, Y. and Yu, L. (2025).
    Advancing Wasserstein convergence analysis of score-based models:
    insights from discretization and second-order acceleration.
    In \textit{Advances in Neural Information Processing Systems 38}.
    \url{https://proceedings.neurips.cc/paper_files/paper/2025/hash/c9d9659d1d960b53e8121469ef1f2df5-Abstract-Conference.html}

    \bibitem{D_Bakry}
    D. Bakry, I. Gentil, and M. Ledoux. Analysis and geometry of Markov diffusion operators. Vol. 348. Grundlehren der Mathematischen Wissenschaften [Fundamental Principles of Mathematical Sciences]. Springer, Cham, 2014, pp. xx+552.
    \url{https://link.springer.com/book/10.1007/978-3-319-00227-9}
    \bibitem{azangulov_deligiannidis_rousseau_2024}
    Azangulov, I., Deligiannidis, G. and Rousseau, J. (2024).
    Convergence of diffusion models under the manifold hypothesis
    in high dimensions.
    arXiv preprint arXiv:2409.18804.
    \url{https://arxiv.org/abs/2409.18804}
    \bibitem{ho_jain_abbeel_2020}
    Ho, J., Jain, A. and Abbeel, P. (2020).
    Denoising diffusion probabilistic models.
    In \textit{Advances in Neural Information Processing Systems 33},
    6840--6851.
    \url{https://proceedings.neurips.cc/paper/2020/hash/4c5bcfec8584af0d967f1ab10179ca4b-Abstract.html}

    \bibitem{li_yan_2024_low_dimensional}
    Li, G. and Yan, Y. (2024).
    Adapting to unknown low-dimensional structures in
    score-based diffusion models.
    In \textit{Advances in Neural Information Processing Systems 37}.
    \url{https://proceedings.neurips.cc/paper_files/paper/2024/hash/e46c303cfd3b93bd32ecf6070318efc3-Abstract-Conference.html}

    \bibitem{pope_et_al_2021}
    Pope, P., Zhu, C., Abdelkader, A., Goldblum, M. and Goldstein, T. (2021).
    The intrinsic dimension of images and its impact on learning.
    In \textit{International Conference on Learning Representations}.
    \url{https://openreview.net/forum?id=XJk19XzGq2J}

    \bibitem{heusel_et_al_2017}
    Heusel, M., Ramsauer, H., Unterthiner, T., Nessler, B. and Hochreiter, S. (2017).
    GANs trained by a two time-scale update rule converge to a local Nash equilibrium.
    In \textit{Advances in Neural Information Processing Systems 30}.
    \url{https://proceedings.neurips.cc/paper/2017/hash/8a1d694707eb0fefe65871369074926d-Abstract.html}
    \bibitem{song_et_al_2021}
    Song, Y., Sohl-Dickstein, J., Kingma, D. P., Kumar, A.,
    Ermon, S. and Poole, B. (2021).
    Score-based generative modeling through stochastic differential equations.
    In \textit{International Conference on Learning Representations}.
    \url{https://openreview.net/forum?id=PxTIG12RRHS}


    \end{thebibliography}
    \end{frame}
  \end{document}
